% ID: SEL-D0007-BB-E04
% ESTADO: validada. S008, 2026-11-07; tres pasadas del mismo autor y revisión matemática, didáctica y visual.
\textbf{a) Derivadas y dominios.}\par\nopagebreak[4]
El argumento del logaritmo debe ser positivo:
\[
\frac{x-1}{x+1}>0\iff x\in(-\infty,-1)\cup(1,\infty).
\]
Sea $u=(x-1)/(x+1)$; entonces $u'=2/(x+1)^2$. Por la regla de la cadena,
\[
f'(x)=\frac{u'}u=\boxed{\frac2{x^2-1}},
\qquad x\in(-\infty,-1)\cup(1,\infty).
\]
No se amplía el dominio a $(-1,1)$ porque la fórmula derivada allí tenga sentido. Tampoco se separa el logaritmo en $\ln(x-1)-\ln(x+1)$ cuando ambos argumentos son negativos.

Para $g$, definida en $\mathbb R$, usamos producto y cadena:
\[
g'(x)=3x^2e^{2x^2}+x^3e^{2x^2}(4x)
=\boxed{x^2(3+4x^2)e^{2x^2}}.
\]

\textbf{b) Parábola.}\par\nopagebreak[4]
Completando el cuadrado,
\[
h(x)=(x+1/2)^2+3/4.
\]
Abre hacia arriba y tiene vértice \boxed{(-1/2,3/4)}, eje $x=-1/2$ y corte vertical \boxed{(0,1)}. No corta el eje horizontal: su discriminante es $1-4=-3<0$, y además $h(x)\geq3/4>0$ en todo $\mathbb R$.
\begin{center}\includegraphics[width=0.60\linewidth]{figuras/s008-parabola.pdf}\end{center}

\textbf{c) Área entre x=-1/2 y x=0.}\par\nopagebreak[4]
Como $h$ es positiva, no hace falta cambiar el signo de la integral:
\[
S=\int_{-1/2}^{0}(x^2+x+1)\,dx
=\left[\frac{x^3}{3}+\frac{x^2}{2}+x\right]_{-1/2}^{0}
=\boxed{\frac5{12}\text{ unidades cuadradas}.}
\]
Comprobación: en ese intervalo $3/4\leq h\leq1$ y su anchura es 1/2, por lo que $3/8\leq S\leq1/2$, coherente con $5/12$.
